\documentclass[UTF8,fontset=fandol,11pt,a4paper]{ctexart}

\usepackage{amsmath,amssymb,mathtools}
\usepackage{booktabs}
\usepackage{geometry}
\usepackage{hyperref}
\usepackage{xcolor}

\geometry{margin=2.35cm}
\hypersetup{colorlinks=true,linkcolor=blue,citecolor=blue,urlcolor=blue}
\newcommand{\ii}{\mathrm{i}}
\newcommand{\cK}{\mathcal K}
\newcommand{\cF}{\mathcal F}
\newcommand{\cG}{\mathcal G}
\newcommand{\cB}{\mathcal B}
\newcommand{\cD}{\mathcal D}
\newcommand{\cN}{\mathcal N}
\newcommand{\cP}{\mathcal P}

\title{连续半波导 DtN 与边界积分耦合的\\导模特征值和误差估计框架}
\author{eig-apost 方法重构稿}
\date{2026年8月11日}

\begin{document}

\maketitle

\begin{abstract}
本文重构二维周期线缺陷波导的主特征值问题。固定准周期参数 $\beta$ 后，首先在左右
半无限周期波导上定义精确 Dirichlet-to-Neumann（DtN）算子，再把它与中心缺陷胞元的
连续边界积分 trace operator 耦合，得到在任何 Fourier 截断、QZ 分解或有限尾端之前
已经成立的算子函数 $\cF(k)$。真实 guided eigenvalue 定义为
$\ker\cF(k)\ne\{0\}$。精确 outgoing Cauchy-data space 是
$\operatorname{graph}\Lambda_\pm(k)$；有限维 trace basis、ordered QZ 和 Riccati
方程只用于近似该图空间，不能反向定义精确 DtN。

离散计算链为：中心与普通胞元的 BIE/Fourier 离散、one-cell generalized pencil、
ordered-QZ stable/unstable deflating subspaces、Dirichlet/Neumann blocks，以及在
Dirichlet chart 安全时形成的 $\Lambda_{\pm,h,M}=N_{s,\pm,M}D_{s,\pm,M}^{-1}$。
chart 不安全时保留增广 Cauchy-data relation 或 Robin realization。本文分别列出中心 BIE、wall
trace 截断、one-cell map、deflating subspace、graph-to-DtN chart、矩阵计算和 root
solve 的误差。由于一阶 DtN 在临界空间 $H^{1/2}\to H^{-1/2}$ 上一般非紧，有限 Fourier
lift 不以完整算子范数逼近它；本文改以 projected DtN consistency、目标解的加权 trace
tail，以及预条件/主符号扣除后的 operator convergence 为理论目标。两层
projected correction 首先只预测 next-level root shift；只有另有 saturation 和
remainder 条件时，才可解释为到连续根的剩余误差。旧的 finite-tail/doubling 主公式
完整保存在 \texttt{legacy-tail.tex}，今后只作 cross-check、reference sequence 或 tail
diagnostic。
\end{abstract}

\section{方法权威、目标和结论等级}

本文使用当前项目局部记号，不继承冻结的 M\"uller--Cauchy 主线记号。Phase 1--2
已经承诺以 DtN+BIE 为首条理论和 estimator 路线，并要求 continuous half-guide DtN、
cell/propagation operator、Riccati/QZ 构造和 discrete DtN matrix 严格分层。本文据此
取代 2026 年 7 月的 finite-tail 主公式。旧稿不删除，其完整归档为
\texttt{legacy-tail.tex}；其中 $N_j=2^j$ 个有限胞元、远端闭合和 Redheffer doubling
以后只提供数值交叉检查，不定义本文的真实谱对象。

本文区分四种结论等级。
\begin{enumerate}
  \item 文献中已对相应 PDE 或有限矩阵对象证明的结果，只按原文假设引用。
  \item 可直接使用的有限维代数恒等式，在本文给出所需条件和坐标约定。
  \item 从连续 BIE 到 discrete QZ/DtN 的收敛、伴随一致性和解析性，是本项目尚需证明的
  命题，不因固定点数值吻合而升级为 theorem。
  \item spectral separation、saturation 和 asymptotic remainder 在没有证明时只能作为
  明确可反驳的数值假设。
\end{enumerate}

本文不组装或运行现有 $A_{\mathrm{def}}$，不执行 locator、DtN wall experiment、
contour count 或 root isolation。下面的 $A_{\mathrm{def}}$ 仅表示未来对 $\cF(k)$ 的
有限维实现所应满足的定义约束。

\section{连续半波导 DtN}

\subsection{几何、trace spaces 和法向约定}

令 $\Omega_0$ 是包含缺陷的一段有界中心区域，左右人工截面分别为
$\Gamma_-$ 和 $\Gamma_+$；$\Omega_-$、$\Omega_+$ 是从相应截面向左、向右延伸的
半无限周期波导。周期方向记为 $y$，周期为 $L_y$，并固定实准周期参数 $\beta$：
\[
  u(x,y+L_y)=e^{\ii\beta L_y}u(x,y).
\]
时间因子取 $e^{-\ii\omega t}$；若 BIE 使用第一类 Hankel 函数，这一选择与 outgoing
Green function 一致。

令 $\nu_-=-e_x$、$\nu_+=e_x$，二者都从中心区域指向相应 half-guide。定义
\[
  X_\pm:=H_\beta^{1/2}(\Gamma_\pm)
  :=\gamma_\pm H_\beta^1(\Omega_\pm),
  \qquad
  Y_\pm:=H_\beta^{-1/2}(\Gamma_\pm)=X_\pm'.
\]
这里 trace space 通过半波导准周期 $H^1$ 空间的迹定义；准周期兼容性由该 trace-space
定义编码，不诉诸临界 $H^{1/2}$ 元素的端点点值。$\partial_{\nu_\pm}$ 始终表示沿
“中心到 half-guide”的定向导数。
Fliss 使用 half-guide 自身的外法向，方向与 $\nu_\pm$ 相反；因此本文与其文献记号的
映射为
\[
  \Lambda_\pm^{\mathrm{Fliss}}=-\Lambda_\pm.
\]
这一负号必须在 theory-to-code map 中显式保存，不能由“左右对称”猜测。

\subsection{projected gap、Wood threshold 和 Dirichlet 例外集}

设 $A_{\mathrm p}(\beta,\theta)$ 是普通二维周期介质在 $y$ 方向取 $\beta$、在 lead
方向取 Bloch 参数 $\theta$ 的胞元算子。固定 $\beta$ 的 bulk projected spectrum 为
\[
  \Sigma_{\mathrm{bulk}}(\beta)
  :=\bigcup_{\theta\in[-\pi/L_x,\pi/L_x]}
    \sigma\!\left(A_{\mathrm p}(\beta,\theta)\right).
\]
projected gap 是 $\mathbb R\setminus\Sigma_{\mathrm{bulk}}(\beta)$ 的开连通分量。
Fliss 的 Proposition 3.1 证明其目标算子的 essential spectrum 正是这一 bulk union，
Theorem 3.5 给出 gap 中 guided eigenfunction 的指数衰减
\cite{Fliss2013}。

projected gap 与 Wood threshold 不是同一个对象。对 BIE 背景 Rayleigh 表示，令
\[
  \alpha_m=\beta+\frac{2\pi m}{L_y},
  \qquad
  \gamma_m(k)^2=k^2-\alpha_m^2.
\]
$\gamma_m(k)=0$ 是所选 background Green/Rayleigh 坐标的 Wood threshold 或分支点；
它可能导致表示失效，但不能用来定义物理 bulk essential spectrum。反之，处在
projected gap 也不自动保证某个准周期 Green evaluator 远离 Wood 点。

另令 $\Sigma_{D,\pm}(\beta)$ 为半波导在 $\Gamma_\pm$ 施加齐次 Dirichlet 数据所得
算子的离散谱。Fliss 的 Theorem 4.1 表明：当 $k^2$ 离开 projected essential spectrum
时，半波导 Dirichlet problem 除一个可数例外集外唯一可解；该例外集正由相应
Dirichlet spectrum 产生。对称胞元是原文给出的特殊无例外情形。Remark 4.2 说明
Robin-to-Robin chart 可避开这些 Dirichlet 例外点\cite{Fliss2013}。

因此实数工作区间 $I$ 必须满足
\[
  k^2\notin
  \Sigma_{\mathrm{bulk}}(\beta)
  \cup\Sigma_{D,-}(\beta)
  \cup\Sigma_{D,+}(\beta),
  \qquad k\in I,
\]
并与所有实际使用的 Wood threshold 保持正距离。做 complex-$k$ 分析时还需选择
一个连通开集 $U\subset\mathbb C$，使 $U\cap\mathbb R\subset I$，所有 square-root
branches 在 $U$ 上由同一 seed 解析延拓，且 half-guide solution operators、中心 BIE
和普通胞元 BIE 在 $U$ 上无 pole。Fliss 直接证明的是实频 well-posedness 和 norm
continuity（Proposition 4.3），不是这里所需的复邻域 holomorphy；后者是本项目必须
补证的命题。

\subsection{精确边值定义和 Cauchy-data graph}

给定 $\phi\in X_\pm$，考虑材料与传输条件均为真实 half-guide 条件的边值问题
\[
  -\Delta u_\pm-k^2q_\pm u_\pm=0\quad\text{in }\Omega_\pm,
  \qquad
  \gamma_\pm u_\pm=\phi,
\]
并要求 $u_\pm$ 属于 projected gap 中的衰减解空间；在 complex-$k$ 邻域则取从该
衰减解空间解析延拓得到的 branch。若此问题唯一可解，定义
\begin{equation}
  \Lambda_\pm(k):X_\pm\longrightarrow Y_\pm,
  \qquad
  \Lambda_\pm(k)\phi
  :=\partial_{\nu_\pm}u_\pm(\phi;k)\big|_{\Gamma_\pm}.
  \label{eq:exact-dtn}
\end{equation}
这一对象只依赖半无限 PDE，不依赖有限 trace basis、Riccati、QZ、BIE 或有限尾端。
Fliss 的 Proposition 4.3 在其相反法向约定下给出
$H_\beta^{1/2}\to H_\beta^{-1/2}$ 的有界性和实参数 norm continuity；Theorem 4.5
给出无界 guided problem 与中心域精确 DtN problem 的双向等价和重数保持
\cite{Fliss2013}。

精确 outgoing/decaying Cauchy-data space 是
\begin{equation}
  \cG_\pm(k)
  :=\operatorname{graph}\Lambda_\pm(k)
  =\{(\phi,\Lambda_\pm(k)\phi):\phi\in X_\pm\}
  \subset X_\pm\times Y_\pm.
  \label{eq:exact-graph}
\end{equation}
式\eqref{eq:exact-graph} 是由式\eqref{eq:exact-dtn} 推出的图表示；它不是用某个有限维
stable trace subspace 反向给出的定义。

\section{连续中心 DtN pencil 与条件化 BIE realization}

\subsection{先定义 exact physical pencil}

令 $V_\beta(\Omega_0)$ 是满足横向准周期条件的中心域能量空间；内部材料界面的传输条件
由中心问题的标准弱形式编码。记 $a_0(k;u,v)$ 为该中心 PDE 的 sesquilinear form，
$\gamma_\pm:V_\beta(\Omega_0)\to X_\pm$ 为完整 wall traces。定义
\begin{equation}
  \langle\cF(k)u,v\rangle
  :=a_0(k;u,v)
  -\sum_{\varsigma\in\{-,+\}}
   \langle\Lambda_\varsigma(k)\gamma_\varsigma u,
   \gamma_\varsigma v\rangle_{\Gamma_\varsigma},
  \qquad
  \cF(k):V_\beta(\Omega_0)\to V_\beta(\Omega_0)'.
  \label{eq:continuous-pencil}
\end{equation}
这里 $\nu_\pm$ 正是中心域的外法向，所以分部积分产生式\eqref{eq:continuous-pencil} 中的
减号；换成 Fliss 的 half-guide 外法向后边界项写成加号。

真实 guided eigenvalue 由 physical pencil 定义：
\begin{equation}
  k_*\in U,
  \qquad
  \ker\cF(k_*)\ne\{0\}.
  \label{eq:true-root}
\end{equation}
Fliss 的 bounded-domain DtN equivalence 在其实参数适用域内为这一 kernel criterion 提供
直接依据；复邻域中的 analytic Fredholm 性仍是本项目证明义务。关键是
式\eqref{eq:continuous-pencil} 在 BIE representation、Fourier truncation、QZ 和矩阵之前
已经成立。

\subsection{continuous BIE trace realization 是待证明的等价表示}

令 $\mathcal X_c$ 是中心 transmission BIE 的连续 density/trace unknown space，
$\mathcal Y_c$ 是其 interface residual space，并令完整 BIE--port residual space 为
\[
  \mathcal Z_c:=\mathcal Y_c\times Y_-\times Y_+.
\]
具体 M\"uller、Calder\'on 或其他无伪频
组合可以改变这些空间及 block 形式，但不能改变物理 trace 接口。记
\[
  \cB(k):\mathcal X_c\to\mathcal Y_c,
  \qquad
  \cD_\pm(k):\mathcal X_c\to X_\pm,
  \qquad
  \cN_\pm(k):\mathcal X_c\to Y_\pm
\]
为 BIE interface residual、完整 wall Dirichlet trace 和沿 $\nu_\pm$ 的 Neumann trace。
记
\[
  \cP_\pm(k):=\cN_\pm(k)-\Lambda_\pm(k)\cD_\pm(k):
  \mathcal X_c\longrightarrow Y_\pm
\]
为连续 port residual。定义条件化的 BIE block schema
\begin{equation}
  \cF_{\mathrm{BIE}}(k):\mathcal X_c\longrightarrow\mathcal Z_c,
  \qquad
  \cF_{\mathrm{BIE}}(k)x=
  \begin{bmatrix}
    \cB(k)x\\
    \cN_-(k)x-\Lambda_-(k)\cD_-(k)x\\
    \cN_+(k)x-\Lambda_+(k)\cD_+(k)x
  \end{bmatrix}.
  \label{eq:continuous-bie-pencil}
\end{equation}
式\eqref{eq:continuous-bie-pencil} 本身还不能取代式\eqref{eq:continuous-pencil} 定义真实根。
本项目必须证明：layer-potential reconstruction 把其 kernel 映到非零物理场；零场不产生
非零 density；每个目标物理解都有 BIE preimage；primal/adjoint pairings 一致。只有这些
kernel--field bridge 条件成立后，才有
\[
  \ker\cF_{\mathrm{BIE}}(k)\simeq\ker\cF(k),
\]
并可把离散 BIE matrix 解释为 physical pencil 的表示。固定矩阵的 raw/reduced Schur
agreement 或 field-energy sanity check 都不能单独证明这一等价性。

\section{从 one-cell BIE 到 discrete DtN}

\subsection{BIE/Fourier 离散和 one-cell generalized pencil}

令 $M$ 是物理 wall trace 的 Fourier order，$K_M=2M+1$。中心和普通周期胞元的 BIE
离散分别还含 boundary quadrature 参数 $h_c,h_{\mathrm{cell}}$。MFS/proxy evaluator
内部的 plane-wave order $M_{\mathrm{pw}}$ 属于 Green evaluator/cell BIE 误差，不是
$M$，也不是 $M_{\mathrm{trace}}$ 或 stable-subspace dimension。

普通胞元的 BIE/Fourier solve 在固定 port normalization 下产生 one-cell scattering
blocks
\[
  \begin{bmatrix}b^L\\a^R\end{bmatrix}
  =
  \begin{bmatrix}R_L&T_{RL}\\T_{LR}&R_R\end{bmatrix}
  \begin{bmatrix}a^L\\b^R\end{bmatrix}.
\]
为消除相位原点的歧义，冻结
\[
  z=\begin{bmatrix}a^L\\b^L\end{bmatrix},
  \qquad a^R=\lambda a^L,
  \qquad b^R=\lambda b^L.
\]
把这两个 Floquet matching relations 代入两条 scattering equations，分别得到
\[
  -R_La^L+b^L=\lambda T_{RL}b^L,
  \qquad
  T_{LR}a^L=\lambda(a^L-R_Rb^L).
\]
因此形成项目当前采用的 generalized pencil candidate
\begin{equation}
  A_{\mathrm{sc}}=
  \begin{bmatrix}-R_L&I\\T_{LR}&0\end{bmatrix},
  \qquad
  B_{\mathrm{sc}}=
  \begin{bmatrix}0&T_{RL}\\I&-R_R\end{bmatrix},
  \qquad
  A_{\mathrm{sc}}z=\lambda B_{\mathrm{sc}}z.
  \label{eq:cell-pencil}
\end{equation}
这里 $\lambda$ 是 lead 方向的 Floquet multiplier，不是目标 guided eigenvalue $k$。
式\eqref{eq:cell-pencil} 是 one-cell scattering relation 的有限维坐标表示；它的
regularity、维数、phase origin 和 normalization 必须作为离散数据记录。

\subsection{ordered QZ 和 stable deflating subspace}

在一个固定 real gap 点，right-decaying multipliers 满足 $|\lambda|<1$，left-decaying
方向对应 positive-$x$ pencil 的 $|\lambda|>1$ cluster。只有当 pencil regular、没有
infinite/indeterminate pair、两个 clusters 各有 $K_M$ 维且与单位圆分离时，才进行
ordered QZ。令
\[
  Z_s=\begin{bmatrix}A_s\\B_s\end{bmatrix},
  \qquad |\lambda|<1,
  \qquad
  Z_u=\begin{bmatrix}A_u\\B_u\end{bmatrix},
  \qquad |\lambda|>1
\]
分别是 right stable 和 positive-$x$ unstable right deflating subspace 的正交基。
需要两次独立 ordering；不能把一次 upper-triangular QZ 的 trailing Schur vectors 自动
当成 invariant subspace。

ordered QZ 是有限矩阵的计算方法。LAPACK 的 xTGSEN 文档明确把 reordered leading
Schur vectors解释为相应 deflating subspaces，并以 generalized Sylvester separation
的 reciprocal estimates 衡量 selected/unselected clusters 的条件性
\cite{LAPACKTGSEN}。Stewart 对 $Az=\lambda Bz$ 给出了 deflating-subspace perturbation
bounds\cite{Stewart1972}。因此本项目的误差门必须记录 cluster separation 或
LAPACK-style \texttt{DIF/sep}；单独记录
$\min_\lambda||\lambda|-1|$ 只是分类 margin，不是一般非正规 pencil 的完整
subspace condition number。

\subsection{Cauchy blocks 和 safe Dirichlet chart}

在当前 amplitude convention 中，wall Dirichlet trace 是 $a+b$。沿中心到 lead 的
定向 Neumann trace在右侧为 $\ii\Gamma(a-b)$，在左侧为
$-\ii\Gamma(a-b)$，其中 $\Gamma=\operatorname{diag}(\gamma_m)$。因此定义
\begin{align}
  D_{s,+}&=A_s+B_s,
  &N_{s,+}&=\ii\Gamma(A_s-B_s),\label{eq:right-cauchy}\\
  D_{s,-}&=A_u+B_u,
  &N_{s,-}&=-\ii\Gamma(A_u-B_u).\label{eq:left-cauchy}
\end{align}
这里 $Z_s$ 在右 half-guide 的中心接口参考面解释；$Z_u$ 是 positive-$x$ unstable
subspace，故向负胞元方向衰减，并须搬运到左 half-guide 的中心接口参考面。若 QZ basis
最初记录在相邻胞元的另一参考面，这一搬运必须作为共同可逆列变换显式记录；它不改变
$\operatorname{range}[D_{s,-};N_{s,-}]$，但漏记其相位会改变 theory-to-code map。左式的
负号来自中心域外法向 $\nu_-=-e_x$，不是由镜像对称性推断。
这一步把 amplitude deflating subspace 转成物理 Cauchy-data subspace。只有
$D_{s,\pm}$ 方阵且 Dirichlet chart 具有可接受的最小奇异值/reciprocal condition
estimate 时，才显式形成
\begin{equation}
  \Lambda_{\pm,h,M}(k)
  =N_{s,\pm,M}(k)D_{s,\pm,M}(k)^{-1},
  \label{eq:discrete-dtn}
\end{equation}
实际实现使用右线性求解，不能形成显式 inverse。若 Schur basis 右乘任意可逆矩阵
$T$，$(D_sT,N_sT)$ 表示同一个 graph，且
$N_sT(D_sT)^{-1}=N_sD_s^{-1}$；因此 DtN 取决于 subspace 和 chart，不取决于
所选 basis。

若两个已对齐的 graph bases 满足
$\widehat D=D+\Delta D$、$\widehat N=N+\Delta N$，则有精确恒等式
\begin{equation}
  \widehat\Lambda-\Lambda
  =(\Delta N-\Lambda\Delta D)(D+\Delta D)^{-1}.
  \label{eq:chart-identity}
\end{equation}
当 $\|D^{-1}\Delta D\|<1$ 时，式\eqref{eq:chart-identity} 给出
\[
  \|\widehat\Lambda-\Lambda\|
  \le
  \frac{(\|\Delta N\|+\|\Lambda\|\,\|\Delta D\|)\|D^{-1}\|}
       {1-\|D^{-1}\Delta D\|}.
\]
这说明 SLP-D 稳定或 deflating residual 很小都不能单独保证 DtN 稳定；Dirichlet
projection 的 transversality 同样是实质条件。

\subsection{unsafe chart 的 graph/Robin realization}

当 $D_{s,\pm}$ 病态或奇异时，不把式\eqref{eq:discrete-dtn} 的巨大数值误差解释成物理
DtN pole，也不把 $\operatorname{range}[D_s;N_s]$ 预先称为某个 operator graph。它首先
只是 discrete Cauchy-data subspace（必要时是相对于 Dirichlet 坐标的 linear relation）。
可保留 coefficient vector $c_\pm$，直接耦合
\begin{equation}
  \cD_{\pm,h}x_h=D_{s,\pm}c_\pm,
  \qquad
  \cN_{\pm,h}x_h=N_{s,\pm}c_\pm.
  \label{eq:aug-graph}
\end{equation}
式\eqref{eq:aug-graph} 是 discrete Cauchy-relation realization。只有 Dirichlet projection
可逆时，它才是式\eqref{eq:discrete-dtn} 的 graph realization。令
$J_M:\mathbb C_D^{K_M}\to\mathbb C_N^{K_M}$ 是冻结的 discrete impedance/Riesz map，负责把
Dirichlet 和 Neumann coefficients 放到相同单位与 Sobolev scaling。也可在
$N_s+\ii\eta J_MD_s$ 安全时使用 Robin chart
\[
  (N_s-\ii\eta J_MD_s)(N_s+\ii\eta J_MD_s)^{-1}.
\]
连续层相应使用固定 isomorphism $J:X_\pm\to Y_\pm$；若不记录 $J_M$ 及其跨层
compatibility，$D_s$ 与 $N_s$ 的直接相加没有坐标无关的含义。
两者都只是式\eqref{eq:exact-graph} 的离散坐标，不是新的 continuous DtN definition。
Fliss 的 Remark 4.2 和后续 RtR 工作为避开 Dirichlet poles 提供了直接先例
\cite{Fliss2013,Fliss2015}。

\subsection{\texorpdfstring{$A_{\mathrm{def}}$}{A-def} 和 block balancing 的位置}

选定中心 density basis、wall Fourier basis、stable Cauchy realization 和 test pairing
以后，先定义未平衡的 discrete BIE--compressed-DtN operator
\begin{equation}
  \cF_{h,M}(k):\mathcal X_{c,h,M}\longrightarrow\mathcal Z_{h,M},
  \label{eq:discrete-f}
\end{equation}
其中
\[
  \mathcal Z_{h,M}:=\mathcal Y_{B,h,M}
  \times\mathbb C_N^{K_M}\times\mathbb C_N^{K_M}
\]
是 center-interface residual 与左右 port residual 的完整 product test space；
$\mathcal Y_{B,h,M}$ 只表示 center BIE residual coefficient space。
它由式\eqref{eq:continuous-bie-pencil} 的 center BIE、wall Fourier restriction 和
式\eqref{eq:discrete-dtn}（或式\eqref{eq:aug-graph}）逐块离散得到。其在所选 trial/test
bases 下的矩阵才记为
\[
  A_{\mathrm{def},h,M}(k).
\]
这才是未来 $A_{\mathrm{def}}$ 的定义：它首先是 $\cF_{\mathrm{BIE}}(k)$ 的离散矩阵；只有
kernel--field bridge 和 regular approximation 成立后，才可称为 physical
$\cF(k)$ 的离散表示。为排除 spectral pollution，除 consistency 外还必须证明 uniform
stability/discrete inf--sup 或相应 resolvent convergence。它不以自身的近奇异性定义新的
连续谱问题。若
$L_{h,M}$、$R_{h,M}$ 可逆，block-balanced matrix
\[
  \widetilde A_{\mathrm{def},h,M}(k)
  =L_{h,M}(k)A_{\mathrm{def},h,M}(k)R_{h,M}(k)
\]
只是 residual 和 unknown 的坐标变换。它不改变 pointwise kernel/root；若 balancing
依赖 $k$ 或 level，则导数、左右向量和 two-level comparison 必须一致地包含这些变换，
不能把 balancing 写入真实 eigenvalue 的定义。

\section{operator-level 误差分解}

\subsection{公共 trace space、投影和临界范数障碍}

令 $Q_M^D:X_\pm\to\mathbb C^{K_M}$ 是 Dirichlet Fourier coefficient map，
$E_M^D:\mathbb C^{K_M}\to X_\pm$ 是 synthesis；$Q_M^N,E_M^N$ 是与
$X_\pm$--$Y_\pm$ duality 相容的 Neumann maps。系数空间分别使用离散
$H^{1/2}$ 和 $H^{-1/2}$ 加权范数。精确 DtN 在该有限 trace space 上的压缩是
\begin{equation}
  \Lambda_{\pm,M}^{\mathrm{ex}}(k)
  :=Q_M^N\Lambda_\pm(k)E_M^D:
  \mathbb C_D^{K_M}\longrightarrow\mathbb C_N^{K_M}.
  \label{eq:dtn-compression}
\end{equation}
离散 half-guide 方法首先必须逼近式\eqref{eq:dtn-compression}，即控制
\begin{equation}
  \varepsilon_{\Lambda,h,M}^{\mathrm{proj}}
  :=\|\Lambda_{\pm,h,M}-\Lambda_{\pm,M}^{\mathrm{ex}}\|_{C^1(\mathfrak K;
  \mathbb C_D^{K_M}\to\mathbb C_N^{K_M})}.
  \label{eq:dtn-projected-target}
\end{equation}

这里必须区分三种 Cauchy-data 对象：无限维物理图
$\cG_\pm(k)=\operatorname{graph}\Lambda_\pm(k)$、retained input 上的 exact compressed
graph
\[
  \cG_{\pm,M}^{\mathrm{ex}}(k)
  :=\operatorname{range}
  \begin{bmatrix}I\\\Lambda_{\pm,M}^{\mathrm{ex}}(k)\end{bmatrix},
\]
以及 QZ/BIE 产生的
$\cG_{\pm,h,M}:=\operatorname{range}[D_{s,\pm};N_{s,\pm}]$。后两者在同一个 weighted
finite coefficient space 中才可比较 symmetric gap。无限维图与有限维 subspace 的
uniform gap 一般不会趋零；$M\to\infty$ 必须使用式\eqref{eq:trace-tail}、strong/Mosco
convergence 或 regular Galerkin approximation。

不能把有限秩 lift $E_M^N\Lambda_{\pm,h,M}Q_M^D$ 在完整
$H^{1/2}\to H^{-1/2}$ operator norm 下与 $\Lambda_\pm$ 作收敛主张。一阶椭圆 DtN 在这对
临界 Sobolev 空间上一般不是紧算子；若这样的有限秩 operator-norm convergence 成立，
极限反而必须紧。归一化高 Fourier 模也直接显示截断余项可保持 $O(1)$。

因此从 projected consistency 回到连续问题还需要以下路线之一。
\begin{enumerate}
  \item 对目标 primal/adjoint eigentraces 证明额外正则性，例如
  $\phi\in H^{1/2+s}$、$\Lambda\phi\in H^{-1/2+s}$，并控制
  \begin{equation}
    \tau_M(\phi;k):=
    \|(I-E_M^DQ_M^D)\phi\|_{H^{1/2}}
    +\|(I-E_M^NQ_M^N)\Lambda(k)\phi\|_{H^{-1/2}}.
    \label{eq:trace-tail}
  \end{equation}
  在上述正则性下，期望 $\tau_M=O(M^{-s})$，但 $s$ 和常数必须由本项目的
  interface/corner regularity 证明或数值估计支持。
  \item 选定并解析保留一个共同的 principal DtN part
  $\Lambda_\pm^{\mathrm{prin}}$，证明余项
  $\Lambda_\pm-\Lambda_\pm^{\mathrm{prin}}$ 在所选空间间紧或更平滑，再只对余项要求
  operator-norm approximation。
  \item 采用 Cauchy graph/strong convergence 和 regular Galerkin spectral
  approximation，而不声称完整 DtN 的有限秩 operator-norm convergence。
\end{enumerate}
若相邻 levels 改变 dimension，仍必须记录 trial prolongation 和 dual-compatible test
restriction；未经 transport 的裸矩阵差不能进入 correction。

\subsection{必须分开的误差层}

在一个紧参数集 $\mathfrak K\Subset U$ 上，建议按以下顺序建立 $C^1(\mathfrak K)$ budget。
\begin{table}[htbp]
  \centering
  \footnotesize
  \caption{连续对象到 root 的误差链。}
  \begin{tabular}{p{0.17\textwidth}p{0.35\textwidth}p{0.37\textwidth}}
    \toprule
    层 & 需要控制的对象 & 可计算诊断与不能替代的证明 \\
    \midrule
    中心 BIE & $\cB,\cD_\pm,\cN_\pm$ 及其 $k$ 导数的离散误差 & boundary refinement、off-surface residual、adjoint test；固定点 agreement 不是 operator convergence \\
    wall trace & Fourier projection/synthesis 和 omitted weighted Cauchy tail & reconstruction、omitted energy、相邻 $M$；制造密度 tail 不是所有解的 operator norm \\
    one-cell & scattering blocks 或 pencil pair $(A_{\mathrm{sc}},B_{\mathrm{sc}})$ 的 $C^1$ 误差 & BIE/proxy/solver refinement、hold-out residual；$M_{\mathrm{pw}}$ 不等于 $M$ \\
    QZ subspace & exact stable/unstable graph 与 computed deflating subspace 的 gap & deflating residual、backward error、\texttt{DIF/sep}、cluster count；unit-circle margin 单独不充分 \\
    DtN chart & projected Cauchy-subspace error 经 $D_s^{-1}$ 传播后的 compressed DtN error & $\sigma_{\min}(D_s)$、式\eqref{eq:chart-identity}；unsafe 时改用 Cauchy relation/RtR \\
    matrix/linear algebra & assembly、evaluation、QZ/solve 和 derivative error & rank policy、solve residual、precision repeat；不得静默换 solver \\
    root solve & 已固定 $\cF_{h,M}$ 的 nonlinear solve error & projected residual、left/right derivative、contour isolation；不等于 operator discretization error \\
    \bottomrule
  \end{tabular}
\end{table}

下面显式冻结 transport，避免把 continuous、compressed 和 discrete maps 混在同一
减法中。BIE trial synthesis 与 center-interface residual restriction 分别写成
\[
  E_{h,M}^X:\mathcal X_{c,h,M}\to\mathcal X_c,
  \qquad
  Q_{B,h,M}:\mathcal Y_c\to\mathcal Y_{B,h,M}
\]
再定义完整 product restriction
\[
  Q_{Z,h,M}:=\operatorname{diag}(Q_{B,h,M},Q_M^N,Q_M^N):
  \mathcal Z_c\longrightarrow\mathcal Z_{h,M}.
\]
exact compressed maps 为
\begin{align*}
  \cB_{h,M}^{\mathrm{ex}}&:=Q_{B,h,M}\cB E_{h,M}^X,\\
  \cD_{\pm,h,M}^{\mathrm{ex}}&:=Q_M^D\cD_\pm E_{h,M}^X,\\
  \cN_{\pm,h,M}^{\mathrm{ex}}&:=Q_M^N\cN_\pm E_{h,M}^X.
\end{align*}
computed maps 均在相同 coefficient spaces 上，故可以定义
\begin{align*}
  E_{\cB}&:=\cB_{h,M}-\cB_{h,M}^{\mathrm{ex}},
  &E_{\cD,\pm}&:=\cD_{\pm,h,M}-\cD_{\pm,h,M}^{\mathrm{ex}},\\
  E_{\cN,\pm}&:=\cN_{\pm,h,M}-\cN_{\pm,h,M}^{\mathrm{ex}},
  &E_{\Lambda,\pm}&:=\Lambda_{\pm,h,M}-\Lambda_{\pm,M}^{\mathrm{ex}}.
\end{align*}
令 projected exact port row 和 computed port row 分别为
\[
  \cP_{\pm,h,M}^{\mathrm{ex}}
  :=\cN_{\pm,h,M}^{\mathrm{ex}}
    -\Lambda_{\pm,M}^{\mathrm{ex}}\cD_{\pm,h,M}^{\mathrm{ex}},
  \qquad
  \cP_{\pm,h,M}:=\cN_{\pm,h,M}
    -\Lambda_{\pm,h,M}\cD_{\pm,h,M}.
\]
则同型对象间有精确恒等式
\begin{align}
  \cP_{\pm,h,M}-\cP_{\pm,h,M}^{\mathrm{ex}}
  &=E_{\cN,\pm}-E_{\Lambda,\pm}\cD_{\pm,h,M}^{\mathrm{ex}}
    -\Lambda_{\pm,h,M}E_{\cD,\pm},\label{eq:port-error}\\
  (\cP_{\pm,h,M}-\cP_{\pm,h,M}^{\mathrm{ex}})'
  &=E_{\cN,\pm}'-E_{\Lambda,\pm}'\cD_{\pm,h,M}^{\mathrm{ex}}
    -E_{\Lambda,\pm}(\cD_{\pm,h,M}^{\mathrm{ex}})'
    -\Lambda_{\pm,h,M}'E_{\cD,\pm}
    -\Lambda_{\pm,h,M}E_{\cD,\pm}'.\label{eq:port-derivative-error}
\end{align}
projected exact port 与 continuous port 的 restriction 之间另有 wall-tail term
\begin{equation}
  \cP_{\pm,h,M}^{\mathrm{ex}}
  -Q_M^N\cP_\pm E_{h,M}^X
  =Q_M^N\Lambda_\pm(I-E_M^DQ_M^D)\cD_\pm E_{h,M}^X.
  \label{eq:port-tail}
\end{equation}
因此 projected BIE block error 是
\[
  \begin{bmatrix}E_{\cB}&E_{\cP,-}&E_{\cP,+}\end{bmatrix}^{\mathsf T},
\]
它是两个从 $\mathcal X_{c,h,M}$ 到 $\mathcal Z_{h,M}$ 的同型 block operators 之差。
若所选 convergence route 另外给出同型的 continuous approximate DtN
$\Lambda_\pm^{\mathrm{app}}:X_\pm\to Y_\pm$，则 physical weak pencil 中纯 DtN 扰动为
\[
  -\sum_\pm\gamma_\pm^*
  (\Lambda_\pm^{\mathrm{app}}-\Lambda_\pm)\gamma_\pm.
\]
其中 $E_{\cP,\pm}:=\cP_{\pm,h,M}-\cP_{\pm,h,M}^{\mathrm{ex}}$。式
\eqref{eq:port-error}--\eqref{eq:port-tail} 是同型空间上的恒等式；各项及其导数的收敛
仍依赖 projected target、wall-tail regularity 和 stability gates。

若 $\varepsilon_{\mathrm{cell}}$ 控制同一个有限维 pencil pair 的 perturbation，
$\operatorname{sep}_{s/u}$ 是 generalized Sylvester separation，而 ordered-QZ backward
error 为 $\varepsilon_{\mathrm{qz}}$，则期望的结构是
\[
  \operatorname{gap}(\cG_{s,h,M}^{\mathrm{exact\ pencil}},
  \cG_{s,h,M}^{\mathrm{computed}})
  \lesssim
  \frac{\varepsilon_{\mathrm{cell}}+\varepsilon_{\mathrm{qz}}}
       {\operatorname{sep}_{s/u}}.
\]
这一形式有 matrix-pencil perturbation theory 支持，但它只控制同一个 finite pencil 的
exact 与 computed deflating subspaces。它不控制 continuous cell relation 到 Fourier
pencil 的截断误差。常数、norm、left/right pairing 和对 $k$ 的一致性仍必须针对
式\eqref{eq:cell-pencil} 证明，当前不能作为已闭合估计。

\subsection{三个主 convergence targets}

第一目标是式\eqref{eq:dtn-projected-target} 趋于零，同时目标 primal/adjoint traces 的
式\eqref{eq:trace-tail} 趋于零。若采用 principal-part 路线，则相应目标改为紧余项的
operator-norm $C^1(\mathfrak K)$ convergence。中心 BIE trace operators 也须在兼容 trial/test
spaces 上一致逼近。利用
\[
  (\Lambda\cD)'=\Lambda'\cD+\Lambda\cD'
\]
所需的第二、第三目标不是未经限定的 full-space finite-rank norm convergence，而是足以
保证目标谱簇 regular approximation 的 primal/adjoint consistency：
\begin{equation}
  \begin{aligned}
    &\Bigl\|Q_{Z,h,M}\cF_{\mathrm{BIE}}(k)x
      -\cF_{h,M}(k)Q_{X,h,M}x\Bigr\|_{\mathcal Z_{h,M}}\to0,\\
    &\Bigl\|Q_{Z,h,M}\cF_{\mathrm{BIE}}'(k)x
      -\cF'_{h,M}(k)Q_{X,h,M}x\Bigr\|_{\mathcal Z_{h,M}}\to0,
  \end{aligned}
  \label{eq:f-targets}
\end{equation}
对正则目标子空间中的 $x\in\mathcal X_c$ 和 $k\in\mathfrak K$ 一致成立，且对相应
adjoint family 也成立；这里
$Q_{X,h,M}:\mathcal X_c\to\mathcal X_{c,h,M}$ 是与 $E_{h,M}^X$ 配套的 trial
restriction，$Q_{Z,h,M}$ 是上面定义的完整 product restriction。还需要
uniform discrete stability、resolvent convergence 或等价的 no-pollution hypothesis。
这些 BIE-level targets 只有与 physical/BIE kernel bridge 合并后，才给出
$\cF_{h,M}$ 对 physical $\cF$ 的 spectral approximation。若能把 physical pencil
$\cF$ 预条件成 identity 加 compact analytic family，则可把目标改写为该 compact part
的 operator-norm/collectively compact convergence。否则必须使用适合非紧 principal part
的 regular Galerkin/Fredholm spectral approximation theorem。

Moskow 的 nonlinear compact-operator result 要求 analytic operator families、collective
compactness、primal/adjoint convergence 和 simple-eigenvalue nondegeneracy
\cite{Moskow2015}。当前 $\cF$ 尚未被证明可预条件成该文的 $I-kT(k)$ compact form，
所以其 Theorem 4.1 只提供证明模板，不是可直接宣告的项目定理。

\section{simple-root correction 和 two-level correction}

本节把 $V_\beta(\Omega_0)'$ 取为连续反对偶，并约定
$\langle f,v\rangle_{V',V}:=f(v)$；$\cF(k)^*:V_\beta(\Omega_0)\to
V_\beta(\Omega_0)'$ 由
\[
  \langle\cF(k)^*y,x\rangle_{V',V}
  :=\overline{\langle\cF(k)x,y\rangle_{V',V}}
\]
定义。设 $k_*$ 是式\eqref{eq:true-root} 的孤立简单根，
$x_*\in\ker\cF(k_*)$、$y_*\in\ker\cF(k_*)^*$，并满足横截条件
\[
  \langle\cF'(k_*)x_*,y_*\rangle_{V',V}\ne0.
\]
对一个小 operator perturbation $\Delta\cF$，simple-root 的一阶位移为
\begin{equation}
  \Delta k
  =-
  \frac{\langle\Delta\cF(k_*)x_*,y_*\rangle_{V',V}}
       {\langle\cF'(k_*)x_*,y_*\rangle_{V',V}}
  +\text{higher-order remainder}.
  \label{eq:operator-correction}
\end{equation}
G\"uttel--Tisseur 的 Theorem 2.20 直接支持 analytic matrix function 的这一有限维
公式\cite{GuettelTisseur2017}；Moskow 给出带额外 compactness、
adjoint convergence 和 multiplicity 假设的 operator-level correction
\cite{Moskow2015}。把二者用于本文仍依赖式\eqref{eq:f-targets}、Fredholm/compact
reduction 或 regular approximation，以及 simple-root persistence。

令 level $\ell$ 汇总 $(h_c,h_{\mathrm{cell}},M)$ 及冻结的 linear algebra policy。
把相邻两层通过已记录的 prolongation/restriction transport 到同一 trial/test space，
所得 operator 记为 $\widehat\cF_\ell$ 和 $\widehat\cF_{\ell+1}$。若 $k_\ell$ 是 coarse
simple root，$x_\ell,y_\ell$ 是相应左右 vectors，则定义
\begin{equation}
  \delta_\ell^{\mathrm{next}}
  :=-
  \frac{\langle
    [\widehat\cF_{\ell+1}(k_\ell)-
     \widehat\cF_\ell(k_\ell)]x_\ell,y_\ell\rangle}
       {\langle\widehat\cF_\ell'(k_\ell)x_\ell,y_\ell\rangle}.
  \label{eq:two-level}
\end{equation}
式\eqref{eq:two-level} 的无附加解释是预测
$k_{\ell+1}-k_\ell$，并应先用实际 next-level matched root 检验。它不是自动估计
$k_*-k_\ell$。

若另有独立依据给出 saturation
\[
  |k_*-k_{\ell+1}|\le q|k_*-k_\ell|,
  \qquad 0\le q<1,
\]
以及 first-order remainder bound
\[
  |(k_{\ell+1}-k_\ell)-\delta_\ell^{\mathrm{next}}|
  \le r_\ell,
\]
才可推出
\begin{equation}
  \frac{\max(0,|\delta_\ell^{\mathrm{next}}|-r_\ell)}{1+q}
  \le |k_*-k_\ell|
  \le
  \frac{|\delta_\ell^{\mathrm{next}}|+r_\ell}{1-q}.
  \label{eq:saturation-bound}
\end{equation}
observed ratios 或有限几个 level 的单调下降不能生成 $q$ 或 $r_\ell$。在它们缺失时，
式\eqref{eq:two-level} 只能称为 next-level correction；只有在独立 reference 上观察到
稳定 $O(1)$ effectivity 并验证渐近区以后，才称为 empirical estimator。

root/linear-algebra error 另行报告。例如在近似根 $\widetilde k_\ell$，
\[
  \delta_\ell^{\mathrm{root}}
  :=-
  \frac{\langle
  \widehat\cF_\ell(\widetilde k_\ell)x_\ell,y_\ell\rangle}
  {\langle
  \widehat\cF_\ell'(\widetilde k_\ell)x_\ell,y_\ell\rangle}
\]
必须明显小于待解释的 discretization correction；否则 two-level effectivity 不可解释。

\section{理论依据和项目证明义务}

\begin{table}[htbp]
  \centering
  \footnotesize
  \caption{来源等级；“直接”不表示已验证项目假设。}
  \begin{tabular}{p{0.18\textwidth}p{0.27\textwidth}p{0.17\textwidth}p{0.25\textwidth}}
    \toprule
    对象 & 来源定位 & 当前等级 & 本项目仍需完成 \\
    \midrule
    projected essential spectrum、half-guide BVP、DtN、中心域等价 & Fliss PDF pp. 7, 10--13：Prop. 3.1, Thm. 4.1, Prop. 4.3, Thm. 4.5 & 原文已有定理 & 映射材料模型、trace spaces、法向；补 complex-$k$ holomorphy \\
    real-gap propagation/Riccati 与 DtN recovery & Fliss PDF pp. 18--19：Thm. 4.12--4.13, Eq. (4.16) & 原文已有定理 & 映射其 FEM cell relation 假设；证明本项目 BIE cell relation 满足相应条件 \\
    absorbing-media Riccati construction & Joly PDF pp. 5, 9, 13, 26；Coatl\'even PDF pp. 3, 8--10\cite{Joly2006,Coatleven2012} & 仅吸收情形有严格结果 & 不把 $\varepsilon=0$ 的 heuristic/conjectural extension 当作 real-gap theorem \\
    ordered QZ deflating subspace & LAPACK xTGSEN；Stewart (1972) & 有限矩阵直接适用 & 证明 pencil regular、cluster dimension 和 $k$-一致 separation；记录 \texttt{DIF/sep} \\
    graph-to-DtN chart & 式\eqref{eq:chart-identity} & 可直接适用的代数命题 & 对齐 graph bases，控制 $D_s^{-1}$ 和 $k$ 导数；unsafe 时使用 graph/RtR \\
    analytic operator spectral approximation & Kato VII Sec. 1.3, Thm. 1.7；Moskow Thm. 3.1/4.1 & 证明模板，不可直接套用 & generalized-pencil reduction、analytic Fredholm/compact reduction、primal/adjoint convergence、root persistence \\
    simple matrix NEP correction & G\"uttel--Tisseur Thm. 2.20 proof & 有限矩阵直接适用 & 保持同一物理 pairing、完整 $F'$、simple root 和 basis transport \\
    saturation / remainder & 当前无项目来源 & 仅数值假设 & 独立证明或保留为 empirical；不得由 observed ratios 生成 \\
    \bottomrule
  \end{tabular}
\end{table}

按当前材料，以下结果必须由本项目证明或以更窄 claim 明确绕开。
\begin{enumerate}
  \item exact half-guide solution operator 和 $\Lambda_\pm(k)$ 在选定复邻域 $U$ 上的
  holomorphy，并区分 physical projected gap、half-guide Dirichlet poles、Wood points、
  cell Dirichlet poles 和 BIE representation poles。
  \item 中心 continuous BIE trace operator 的 Fredholm property、representation
  injectivity、kernel--field equivalence 和 adjoint consistency。
  \item one-cell BIE/Fourier pencil pair及其 $k$ 导数向 continuous cell relation 的一致
  收敛；ordered-QZ stable/unstable clusters 在 $U$ 上维数固定且 uniformly separated。
  \item deflating-subspace error 经 Cauchy transform 和 safe chart 传播到
  式\eqref{eq:dtn-projected-target} 的定量界；同时以式\eqref{eq:trace-tail}、principal-part
  compact remainder 或 regular Galerkin theory 连接到 $M\to\infty$。
  \item 中心、trace、cell 和 projected DtN 在公共 trial/test space 上给出
  $\cF,\cF'$ 的 primal/adjoint regular approximation；禁止复活完整
  $H^{1/2}\to H^{-1/2}$ finite-rank operator-norm target。
  \item simple-root correction 的 remainder；若要把 two-level correction 升级为
  remaining-error estimator，还需 saturation 或独立 asymptotic tail model。
\end{enumerate}

\section{finite-tail 方法的新位置}

归档的 \texttt{legacy-tail.tex} 保留 finite scattering blocks、远端 Dirichlet/Robin
闭合、Redheffer doubling、augmented finite-tail NEP、complex locator/root protocol 和
finite-tail projected correction 的全部历史推导。它在新路线中只有三种用途：
\begin{enumerate}
  \item 与 ordered-QZ/Riccati discrete half-guide graph 作 same-cell infinity-treatment
  cross-check；共享 one-cell BIE 时不算独立 physical reference。
  \item 作为有限尾长 $N$ 的 reference sequence，检查由 stable multiplier 预测的 decay
  和 termination sensitivity。
  \item 在真实 solved density 上作可选 tail diagnostic，帮助识别 close-to-band-edge
  或 weak-separation 情形。
\end{enumerate}
它不再定义 exact-half-guide lift，不再把 $k_{\infty,h}$ 当作连续谱对象，也不再以
finite-tail level difference 定义主 estimator。

\section{恢复数值 I4 前的最小理论门}

新框架已经足以重新设计 $A_{\mathrm{def}}$ 的 block roles，但尚未授权组装或 locator。
恢复 I4 前至少需要关闭以下 formulation blockers。
\begin{enumerate}
  \item 冻结式\eqref{eq:continuous-pencil} 的具体中心 BIE spaces、wall trace rows、
  source/target normal signs 和 Fliss/code sign map；先保持 physical
  式\eqref{eq:continuous-pencil} 为真实问题，再给出 BIE kernel--field equivalence 的证明
  或明确的数值论文级替代门。
  \item 写出 package/test one-cell scattering blocks 到式\eqref{eq:cell-pencil}、
  式\eqref{eq:right-cauchy}--\eqref{eq:left-cauchy} 的逐块 theory-to-code map；预注册
  regularity、stable/unstable count、\texttt{DIF/sep} 和 graph residual gates。
  \item 对每个 $k$ 选择 safe DtN chart；否则改用增广 Cauchy relation 或 Robin
  realization。冻结 impedance map $J_M$，不能在 Dirichlet chart 失效后继续形成
  $N_sD_s^{-1}$。
  \item 冻结 wall trace common space、prolongation/dual restriction、$M$ 与
  $M_{\mathrm{pw}}$ 的不同角色，以及 $\cF'$ 中所有 $k$ 依赖。
  \item 选择并证明 projected consistency + eigentrace regularity、principal-part
  subtraction 或 regular Galerkin 三条 convergence route 之一；证明在该意义下
  $\cF_{h,M}$ 是 $\cF$ 的一致离散，且 block balancing 仅为坐标变换。在此以前不解释
  任何 $\sigma_{\min}$ dip。
\end{enumerate}

全部门关闭以后，I4 才能按“构造一个由式\eqref{eq:continuous-pencil} 派生的 balanced
matrix evaluator，再做 bounded locator 和 analytic-disk readiness”的顺序恢复。
operator convergence、saturation 和 certified remainder 可以晚于第一个真实 root，
但若它们未闭合，后续结果只能是 qualified discrete root 和 empirical correction，
不能声称 theorem-level estimator 或 certified interval。

\begin{thebibliography}{9}

\bibitem{Fliss2013}
S.~Fliss,
\newblock A Dirichlet-to-Neumann approach for the exact computation of guided modes in photonic crystal waveguides,
\newblock \emph{SIAM Journal on Scientific Computing}, 35(2), B438--B461, 2013.
\newblock DOI: \href{https://doi.org/10.1137/12086697X}{10.1137/12086697X}.

\bibitem{Fliss2015}
S.~Fliss, D.~Klindworth, and K.~Schmidt,
\newblock Robin-to-Robin transparent boundary conditions for the computation of guided modes in photonic crystal wave-guides,
\newblock \emph{BIT Numerical Mathematics}, 55, 81--115, 2015.
\newblock DOI: \href{https://doi.org/10.1007/s10543-014-0521-1}{10.1007/s10543-014-0521-1}.

\bibitem{Joly2006}
P.~Joly, J.-R.~Li, and S.~Fliss,
\newblock Exact boundary conditions for periodic waveguides containing a local perturbation,
\newblock \emph{Communications in Computational Physics}, 1(6), 945--973, 2006.
\newblock DOI: \href{https://doi.org/10.4208/cicp.2006.v1.p945}{10.4208/cicp.2006.v1.p945}.

\bibitem{Coatleven2012}
J.~Coatl\'even,
\newblock Helmholtz equation in periodic media with a line defect,
\newblock \emph{Journal of Computational Physics}, 231(4), 1675--1704, 2012.
\newblock DOI: \href{https://doi.org/10.1016/j.jcp.2011.10.022}{10.1016/j.jcp.2011.10.022}.

\bibitem{Stewart1972}
G.~W.~Stewart,
\newblock On the sensitivity of the eigenvalue problem $Ax=\lambda Bx$,
\newblock \emph{SIAM Journal on Numerical Analysis}, 9(4), 669--686, 1972.
\newblock DOI: \href{https://doi.org/10.1137/0709056}{10.1137/0709056}.

\bibitem{LAPACKTGSEN}
E.~Anderson et al.,
\newblock \emph{LAPACK Users' Guide}, 3rd ed., SIAM, 1999; xTGSEN documentation,
\newblock \href{https://www.netlib.org/lapack/lug/node58.html}{Deflating subspaces and condition numbers}.

\bibitem{Kato1995}
T.~Kato,
\newblock \emph{Perturbation Theory for Linear Operators}, 2nd ed., Springer, 1995.
\newblock DOI: \href{https://doi.org/10.1007/978-3-642-66282-9}{10.1007/978-3-642-66282-9}.

\bibitem{GuettelTisseur2017}
S.~G\"uttel and F.~Tisseur,
\newblock The nonlinear eigenvalue problem,
\newblock \emph{Acta Numerica}, 26, 1--94, 2017.
\newblock DOI: \href{https://doi.org/10.1017/S0962492917000034}{10.1017/S0962492917000034}.

\bibitem{Moskow2015}
S.~Moskow,
\newblock Nonlinear eigenvalue approximation for compact operators,
\newblock \emph{Journal of Mathematical Physics}, 56, 113512, 2015.
\newblock DOI: \href{https://doi.org/10.1063/1.4936304}{10.1063/1.4936304}.

\end{thebibliography}

\end{document}
